\documentclass[11pt]{article}

\usepackage[a4paper,margin=1in]{geometry}
\usepackage{amsmath,amssymb,amsthm,mathtools,bm}
\usepackage[T1]{fontenc}
\usepackage{lmodern}
\usepackage{microtype}
\usepackage{hyperref}

\allowdisplaybreaks
\numberwithin{equation}{section}

\newtheorem{theorem}{Theorem}[section]
\newtheorem{proposition}[theorem]{Proposition}
\newtheorem{lemma}[theorem]{Lemma}
\newtheorem{corollary}[theorem]{Corollary}
\newtheorem{remark}[theorem]{Remark}
\newtheorem{definition}[theorem]{Definition}

\newcommand{\R}{\mathbb R}
\newcommand{\C}{\mathbb C}
\newcommand{\N}{\mathbb N}
\newcommand{\A}{\mathbf A}
\newcommand{\I}{\mathbf I}
\newcommand{\Bmat}{\mathbf B}
\newcommand{\Hmat}{\mathbf H}
\newcommand{\Tmat}{\mathbf T}
\newcommand{\PiM}{\mathbf \Pi}
\newcommand{\tridiag}{\operatorname{tridiag}}
\newcommand{\diag}{\operatorname{diag}}
\newcommand{\smin}{s_{\min}}

\title{Simply Connected Pseudospectral Components and\\
the Connectedness Threshold Problem for the\\
Open-Boundary Hatano--Nelson Matrix}
\author{}
\date{}

\begin{document}
\maketitle

\begin{abstract}
For the real nonnormal tridiagonal Toeplitz matrix
\[
\A_n(a)=\tridiag(a,0,1),\qquad n>1,\quad 0<a<1,
\]
we prove that every connected component of every pseudospectrum
\[
\Sigma_\varepsilon^{(n)}
=
\sigma_\varepsilon\!\bigl(\A_n(a)\bigr)
\]
is simply connected; in fact, every component is contractible.  The key analytic
statement is the vertical monotonicity of the least singular value:
for every fixed \(x\in\R\),
\[
y\longmapsto \smin\!\bigl((x+iy)\I_n-\A_n(a)\bigr)
\quad\text{is nondecreasing on }[0,\infty).
\]
The proof combines a topological reduction, a singular-vector current identity,
a Perron--Frobenius analysis on the imaginary axis, continuation of phase
inequalities in the \(x\)-variable, and a no-zero theorem for ground-state
eigenvectors obtained from one-sided Schur complements of
\[
\Hmat_x(y)=\bigl((x+iy)\I_n-\A_n(a)\bigr)^*
\bigl((x+iy)\I_n-\A_n(a)\bigr).
\]

We also merge into the same framework a reduction of the
\(\varepsilon\)-connectedness threshold problem to a one-dimensional real-axis
barrier problem.  More precisely, if
\[
I_\varepsilon^{(n)}
=
\Sigma_\varepsilon^{(n)}\cap \R
=
\left\{x\in\R:\ \smin\!\bigl(x\I_n-\A_n(a)\bigr)<\varepsilon\right\},
\]
then \(\Sigma_\varepsilon^{(n)}\) is connected if and only if \(I_\varepsilon^{(n)}\)
is connected, and hence if and only if \(\varepsilon\) exceeds an explicit maximal
barrier height on the real axis.  Finally, we record a factorization of the real-axis
boundary equation into two ordinary characteristic polynomials, together with
structured pentadiagonal determinant formulas useful for further algebraic study.
\end{abstract}

\section{Introduction and main results}

Fix throughout \(0<a<1\) and \(n\ge2\), and let
\[
\A_n(a)=
\begin{bmatrix}
0 & 1 \\
a & 0 & \ddots \\
& \ddots & \ddots & \ddots \\
&& \ddots & 0 & 1 \\
&&& a & 0
\end{bmatrix}\in\R^{n\times n}.
\]
For \(z=x+iy\in\C\), define
\[
s_n(x,y)=\smin\!\bigl((x+iy)\I_n-\A_n(a)\bigr),
\qquad
s_n(x)=s_n(x,0),
\]
and let
\[
\Sigma_\varepsilon^{(n)}
=
\sigma_\varepsilon\!\bigl(\A_n(a)\bigr)
=
\{z\in\C:\ s_n(\Re z,\Im z)<\varepsilon\}.
\]
When \(n\) and \(a\) are fixed, we write
\[
\A=\A_n(a),\qquad
s(x,y)=s_n(x,y),\qquad
\lambda_1(x,y)=s(x,y)^2,\qquad
\Sigma_\varepsilon=\Sigma_\varepsilon^{(n)}.
\]

Our first main theorem is topological.

\begin{theorem}[Simply connected components]\label{thm:main}
For every \(\varepsilon>0\), every connected component of
\[
\Sigma_\varepsilon=\sigma_\varepsilon\!\bigl(\A_n(a)\bigr)
\]
is simply connected.  In fact, every connected component is contractible.
\end{theorem}

The proof is reduced to the following analytic monotonicity statement.

\begin{theorem}[Vertical monotonicity]\label{thm:vertical}
For every fixed \(x\in\R\), the function
\[
[0,\infty)\ni y\longmapsto s(x,y)
\]
is nondecreasing.  Equivalently,
\[
\partial_y\lambda_1(x,y)\ge0,
\qquad x\in\R,\ y>0.
\]
\end{theorem}

The connectedness-threshold problem is likewise reduced to a one-dimensional
question.  Define the real slice
\[
I_\varepsilon^{(n)}
:=
\Sigma_\varepsilon^{(n)}\cap\R
=
\{x\in\R:\ s_n(x)<\varepsilon\}.
\]
Write the real eigenvalues of \(\A_n(a)\) as
\[
\mu_k^{(n)}=2\sqrt a\,\cos\!\left(\frac{k\pi}{n+1}\right),
\qquad
k=1,\dots,n,
\]
so that
\[
\mu_1^{(n)}>\mu_2^{(n)}>\cdots>\mu_n^{(n)}.
\]
For \(k=1,\dots,n-1\), define the \(k\)-th real-axis barrier
\[
\beta_{n,k}(a)
:=
\max_{x\in[\mu_{k+1}^{(n)},\,\mu_k^{(n)}]} s_n(x),
\]
and let
\[
\eta_n(a):=\max_{1\le k\le n-1}\beta_{n,k}(a).
\]

The results proved below imply:

\begin{proposition}[Barrier criterion]\label{prop:barrier-intro}
Assume Theorem~\ref{thm:vertical}. Then, for every \(n\ge2\) and
\(\varepsilon>0\),
\[
\Sigma_\varepsilon^{(n)}\ \text{is connected}
\quad\Longleftrightarrow\quad
I_\varepsilon^{(n)}\ \text{is connected}
\quad\Longleftrightarrow\quad
\varepsilon>\eta_n(a).
\]
\end{proposition}

Thus the unresolved part of the threshold problem is the monotonicity in \(n\) of
\(\eta_n(a)\).

\begin{corollary}[Reduction of the threshold problem]\label{cor:threshold-intro}
If
\[
\eta_{n+1}(a)\le \eta_n(a)\qquad (n\ge2),
\]
then for every \(\varepsilon>0\) the set
\[
\mathcal C_{a,\varepsilon}
=
\{n\ge2:\ \Sigma_\varepsilon^{(n)}\ \text{is connected}\}
\]
is an upper tail in \(\N\). More precisely,
\[
\Sigma_\varepsilon^{(n)}\ \text{is connected}
\iff
\eta_n(a)<\varepsilon,
\]
so
\[
N_c(a,\varepsilon):=\min\{n\ge2:\ \eta_n(a)<\varepsilon\}
\]
would satisfy
\[
n\ge N_c(a,\varepsilon)
\quad\Longrightarrow\quad
\Sigma_\varepsilon^{(n)}\ \text{is connected}.
\]
\end{corollary}

Sections~\ref{sec:topology}--\ref{sec:proof-vertical} establish
Theorems~\ref{thm:vertical} and \ref{thm:main}, and also prove the
barrier criterion and the threshold reduction.  Appendices~\ref{app:factor},
\ref{app:pentadiagonal}, and \ref{app:even-central} collect the real-axis
factorization and determinant identities.

\section{Spectrum, symmetry, and topological consequences of vertical monotonicity}
\label{sec:topology}

Throughout this section we fix \(n\ge2\) and \(0<a<1\), and abbreviate
\[
\A=\A_n(a),\qquad
s(x,y)=\smin\!\bigl((x+iy)\I_n-\A\bigr),\qquad
\Sigma_\varepsilon=\{x+iy:\ s(x,y)<\varepsilon\}.
\]
We first record the spectrum.

\begin{lemma}\label{lem:spectrum}
Let
\[
\Tmat_a=
\diag\!\bigl(1,a^{1/2},a,\dots,a^{(n-1)/2}\bigr).
\]
Then
\[
\Tmat_a^{-1}\A_n(a)\Tmat_a
=
\sqrt a\,
\tridiag(1,0,1).
\]
Hence the eigenvalues of \(\A_n(a)\) are real and simple:
\[
\sigma\!\bigl(\A_n(a)\bigr)
=
\left\{
2\sqrt a\,\cos\!\left(\frac{k\pi}{n+1}\right):
k=1,\dots,n
\right\}.
\]
\end{lemma}

\begin{proof}
A direct entrywise computation gives
\[
(\Tmat_a^{-1}\A_n(a)\Tmat_a)_{j,j+1}
=
a^{-(j-1)/2}a^{j/2}=\sqrt a
\]
and
\[
(\Tmat_a^{-1}\A_n(a)\Tmat_a)_{j+1,j}
=
a^{-j/2}\,a\,a^{(j-1)/2}=\sqrt a.
\]
Thus
\[
\Tmat_a^{-1}\A_n(a)\Tmat_a
=
\sqrt a\,\tridiag(1,0,1).
\]
The Dirichlet eigenvalues of \(\tridiag(1,0,1)\) are
\(
2\cos\!\bigl(k\pi/(n+1)\bigr)
\),
\(k=1,\dots,n\), and they are simple.
\end{proof}

\begin{lemma}[Complex conjugation symmetry]\label{lem:conj}
For all \(x,y\in\R\),
\[
s(x,-y)=s(x,y).
\]
Hence \(\Sigma_\varepsilon\) is invariant under complex conjugation.
\end{lemma}

\begin{proof}
Since \(\A\) is real,
\[
\overline{(x+iy)\I_n-\A}=(x-iy)\I_n-\A,
\]
and complex conjugation preserves singular values.
\end{proof}

\begin{lemma}[Every connected component contains an eigenvalue]
\label{lem:component-eigenvalue}
Every connected component of \(\Sigma_\varepsilon\) contains an eigenvalue of \(\A\).
In particular, every connected component of \(\Sigma_\varepsilon\) meets the real axis.
\end{lemma}

\begin{proof}
Since
\[
\smin(z\I_n-\A)\ge |z|-\|\A\|_2,
\]
the set \(\Sigma_\varepsilon\) is bounded. Let \(\Omega\) be a connected component of
\(\Sigma_\varepsilon\), and suppose that
\[
\Omega\cap\sigma(\A)=\varnothing.
\]
Then
\[
\varphi(z):=\|(z\I_n-\A)^{-1}\|_2
\]
is continuous on \(\overline\Omega\), subharmonic on \(\Omega\), and satisfies
\[
\varphi(z)>\varepsilon^{-1},
\qquad z\in\Omega,
\]
because \(z\in\Sigma_\varepsilon\) is equivalent to \(\smin(z\I_n-\A)<\varepsilon\).

Since \(\Omega\) is a connected component of the open set \(\Sigma_\varepsilon\),
\[
\partial\Omega\subset\partial\Sigma_\varepsilon.
\]
By continuity of the least singular value,
\[
\smin(z\I_n-\A)=\varepsilon,
\qquad z\in\partial\Sigma_\varepsilon,
\]
and therefore
\[
\varphi(z)=\varepsilon^{-1},
\qquad z\in\partial\Omega.
\]
This contradicts the maximum principle for subharmonic functions on bounded
domains.  Hence \(\Omega\) contains an eigenvalue of \(\A\), and by
Lemma~\ref{lem:spectrum} that eigenvalue is real.
\end{proof}

We now state the first main topological consequence of
Theorem~\ref{thm:vertical}.

\begin{proposition}[Vertical sections and contractibility]
\label{prop:topological-reduction}
Assume Theorem~\ref{thm:vertical}.  Then for every \(x\in\R\), the vertical section
\[
\Sigma_\varepsilon\cap(\{x\}+i\R)
\]
is either empty or of the form
\[
\{x+iy:\ |y|<h_\varepsilon(x)\}
\]
for some \(h_\varepsilon(x)>0\).

Consequently, every connected component \(\Omega\) of \(\Sigma_\varepsilon\) is of the
form
\[
\Omega
=
\{x+iy:\ x\in I_\Omega,\ |y|<h_\Omega(x)\},
\]
where \(I_\Omega\subset\R\) is an open interval and \(h_\Omega:I_\Omega\to(0,\infty)\).
In particular, \(\Omega\) is contractible.
\end{proposition}

\begin{proof}
Fix \(x\in\R\) and let
\[
V_x:=\{y\in\R:\ x+iy\in\Sigma_\varepsilon\}.
\]
By Lemma~\ref{lem:conj}, the set \(V_x\) is symmetric under \(y\mapsto -y\).
By Theorem~\ref{thm:vertical}, the set
\[
\{y\ge0:\ s(x,y)<\varepsilon\}
\]
is either empty or an interval \([0,h)\).  Therefore \(V_x\) is either empty or
\(
(-h,h)
\)
for some \(h>0\), which proves the first claim.

Now let \(\Omega\) be a connected component of \(\Sigma_\varepsilon\), and define
\[
I_\Omega
=
\{x\in\R:\ \Omega\cap(\{x\}+i\R)\neq\varnothing\}.
\]
The projection \(x+iy\mapsto x\) is continuous, so \(I_\Omega\) is connected.
Since \(\Omega\) is open, \(I_\Omega\) is open.  Hence \(I_\Omega\) is an open interval.

By Lemma~\ref{lem:component-eigenvalue}, \(\Omega\) contains a real eigenvalue
\(x_0\).  Since complex conjugation preserves \(\Sigma_\varepsilon\), the set
\(\overline{\Omega}\) is also a connected component of \(\Sigma_\varepsilon\); because
\(\overline{\Omega}\) contains the same real point \(x_0\), we must have
\[
\overline{\Omega}=\Omega.
\]
Thus each fiber
\[
\Omega_x:=\{y\in\R:\ x+iy\in\Omega\},\qquad x\in I_\Omega,
\]
is symmetric with respect to \(0\).

For fixed \(x\in I_\Omega\), the full section \(V_x\) is connected, and it is partitioned
into the disjoint open slices belonging to the connected components of
\(\Sigma_\varepsilon\).  Since \(\Omega_x\) is a nonempty open symmetric subset of the
interval \(V_x\), it must coincide with all of \(V_x\).  Hence
\[
\Omega_x=(-h_\Omega(x),h_\Omega(x))
\]
for some \(h_\Omega(x)>0\), giving the stated representation.

Finally,
\[
H:[0,1]\times\Omega\to\Omega,\qquad
H(t,x+iy)=x+i(1-t)y,
\]
is a deformation retraction of \(\Omega\) onto \(\Omega\cap\R=I_\Omega\).  Therefore
\(\Omega\) is contractible.
\end{proof}

We next reduce connectedness of the full pseudospectrum to connectedness of the
real slice.

\begin{lemma}[Connectedness of the full pseudospectrum versus the real slice]
\label{lem:real-slice-connected}
Assume Theorem~\ref{thm:vertical}. Define
\[
I_\varepsilon^{(n)}
:=
\Sigma_\varepsilon^{(n)}\cap\R
=
\{x\in\R:\ s_n(x)<\varepsilon\}.
\]
Then
\[
\Sigma_\varepsilon^{(n)}\ \text{is connected}
\quad\Longleftrightarrow\quad
I_\varepsilon^{(n)}\ \text{is connected}.
\]
\end{lemma}

\begin{proof}
Let \(\pi:\C\to\R\) be the projection \(\pi(x+iy)=x\).  We first show that
\[
\pi(\Sigma_\varepsilon^{(n)})=I_\varepsilon^{(n)}.
\]
If \(z=x+iy\in\Sigma_\varepsilon^{(n)}\), then \(s_n(x,y)<\varepsilon\).  By
Lemma~\ref{lem:conj} and Theorem~\ref{thm:vertical},
\[
s_n(x,0)\le s_n(x,|y|)=s_n(x,y)<\varepsilon,
\]
so \(x\in I_\varepsilon^{(n)}\).  Hence
\[
\pi(\Sigma_\varepsilon^{(n)})\subset I_\varepsilon^{(n)}.
\]
Conversely, if \(x\in I_\varepsilon^{(n)}\), then \(x+0i\in\Sigma_\varepsilon^{(n)}\), so
\(x\in \pi(\Sigma_\varepsilon^{(n)})\).  Therefore
\[
\pi(\Sigma_\varepsilon^{(n)})=I_\varepsilon^{(n)}.
\]

If \(\Sigma_\varepsilon^{(n)}\) is connected, then
\(I_\varepsilon^{(n)}\) is connected as the continuous image of a connected set.

Conversely, suppose that \(I_\varepsilon^{(n)}\) is connected.  Since it is a connected
subset of \(\R\), it is an interval.  Let
\[
z_0=x_0+iy_0,\qquad z_1=x_1+iy_1
\]
be any two points of \(\Sigma_\varepsilon^{(n)}\).  By the first part of the proof,
\(x_0,x_1\in I_\varepsilon^{(n)}\).

For \(\tau\in[0,1]\), Lemma~\ref{lem:conj} and
Theorem~\ref{thm:vertical} imply
\[
s_n(x_0,\tau y_0)=s_n(x_0,|\tau y_0|)
\le s_n(x_0,|y_0|)<\varepsilon,
\]
and similarly
\[
s_n(x_1,\tau y_1)=s_n(x_1,|\tau y_1|)
\le s_n(x_1,|y_1|)<\varepsilon.
\]
Thus the vertical segments from \(z_0\) to \(x_0\) and from \(x_1\) to \(z_1\) lie in
\(\Sigma_\varepsilon^{(n)}\).

Since \(I_\varepsilon^{(n)}\) is an interval, the horizontal segment
\[
\{(1-\tau)x_0+\tau x_1:\ \tau\in[0,1]\}
\]
lies in \(I_\varepsilon^{(n)}\subset \Sigma_\varepsilon^{(n)}\).

Concatenating these three paths gives a path in \(\Sigma_\varepsilon^{(n)}\) joining
\(z_0\) to \(z_1\).  Hence \(\Sigma_\varepsilon^{(n)}\) is path-connected, and therefore
connected.
\end{proof}

\begin{lemma}[Every component of the real slice contains an eigenvalue]
\label{lem:real-component-eigenvalue}
Assume Theorem~\ref{thm:vertical}. Then every connected component of
\[
I_\varepsilon^{(n)}
=
\Sigma_\varepsilon^{(n)}\cap\R
\]
contains an eigenvalue of \(\A_n(a)\).
\end{lemma}

\begin{proof}
Let \(J\) be a connected component of \(I_\varepsilon^{(n)}\), and set
\[
\Omega_J:=\Sigma_\varepsilon^{(n)}\cap \pi^{-1}(J).
\]
By the vertical-segment argument in the proof of
Lemma~\ref{lem:real-slice-connected}, \(\Omega_J\) is path-connected: every point of
\(\Omega_J\) can be joined vertically to the real axis, and any two real points in \(J\)
can be joined inside \(J\).

Moreover,
\[
\Omega_J=(\pi|_{\Sigma_\varepsilon^{(n)}})^{-1}(J).
\]
Since \(J\) is open and closed in \(I_\varepsilon^{(n)}\), it follows that \(\Omega_J\) is open
and closed in \(\Sigma_\varepsilon^{(n)}\).  Hence \(\Omega_J\) is a connected component of
\(\Sigma_\varepsilon^{(n)}\).

By Lemma~\ref{lem:component-eigenvalue}, \(\Omega_J\) contains an eigenvalue of
\(\A_n(a)\).  By Lemma~\ref{lem:spectrum}, all eigenvalues of \(\A_n(a)\) are real, so
this eigenvalue lies in
\[
\Omega_J\cap\R = J.
\]
\end{proof}

We can now formulate the exact real-axis barrier criterion.

\begin{proposition}[Exact barrier criterion]\label{prop:barrier}
Assume Theorem~\ref{thm:vertical}.  For \(k=1,\dots,n-1\), define
\[
\mu_k^{(n)}=2\sqrt a\,\cos\!\left(\frac{k\pi}{n+1}\right),
\]
\[
\beta_{n,k}(a)
=
\max_{x\in[\mu_{k+1}^{(n)},\,\mu_k^{(n)}]} s_n(x),
\qquad
\eta_n(a):=\max_{1\le k\le n-1}\beta_{n,k}(a).
\]
Then, for every \(\varepsilon>0\),
\[
I_\varepsilon^{(n)}\ \text{is connected}
\quad\Longleftrightarrow\quad
\varepsilon>\eta_n(a).
\]
Equivalently,
\[
\Sigma_\varepsilon^{(n)}\ \text{is connected}
\quad\Longleftrightarrow\quad
\varepsilon>\eta_n(a).
\]
\end{proposition}

\begin{proof}
By Lemma~\ref{lem:real-slice-connected}, it is enough to work with
\(I_\varepsilon^{(n)}\).

Assume first that \(\varepsilon>\eta_n(a)\).  Then for every \(k=1,\dots,n-1\),
\[
s_n(x)\le \beta_{n,k}(a)<\varepsilon,
\qquad
x\in[\mu_{k+1}^{(n)},\,\mu_k^{(n)}].
\]
Hence
\[
[\mu_{k+1}^{(n)},\,\mu_k^{(n)}]\subset I_\varepsilon^{(n)},
\qquad k=1,\dots,n-1,
\]
so
\[
[\mu_n^{(n)},\,\mu_1^{(n)}]\subset I_\varepsilon^{(n)}.
\]
Every connected component of \(I_\varepsilon^{(n)}\) contains an eigenvalue of
\(\A_n(a)\) by Lemma~\ref{lem:real-component-eigenvalue}.  Since all eigenvalues lie
in the connected subset \([\mu_n^{(n)},\mu_1^{(n)}]\subset I_\varepsilon^{(n)}\), there can
be only one connected component. Thus \(I_\varepsilon^{(n)}\) is connected.

Conversely, assume that \(\varepsilon\le \eta_n(a)\).  Choose
\(k\in\{1,\dots,n-1\}\) such that
\[
\beta_{n,k}(a)=\eta_n(a).
\]
Then there exists \(x_*\in[\mu_{k+1}^{(n)},\mu_k^{(n)}]\) with
\[
s_n(x_*)=\beta_{n,k}(a)\ge \varepsilon,
\]
so \(x_*\notin I_\varepsilon^{(n)}\).

On the other hand, since \(\mu_k^{(n)}\) and \(\mu_{k+1}^{(n)}\) are eigenvalues of
\(\A_n(a)\),
\[
s_n(\mu_k^{(n)})=s_n(\mu_{k+1}^{(n)})=0<\varepsilon,
\]
and therefore
\[
\mu_k^{(n)},\ \mu_{k+1}^{(n)}\in I_\varepsilon^{(n)}.
\]
If \(I_\varepsilon^{(n)}\) were connected, then it would be an interval and hence would
contain all of
\(
[\mu_{k+1}^{(n)},\mu_k^{(n)}]
\),
in particular \(x_*\), a contradiction.  Therefore \(I_\varepsilon^{(n)}\) is
disconnected.
\end{proof}

\begin{corollary}[Reduction of the threshold problem]\label{cor:threshold}
If
\[
\eta_{n+1}(a)\le \eta_n(a)\qquad(n\ge2),
\]
then for every \(\varepsilon>0\) the set
\[
\mathcal C_{a,\varepsilon}
=
\{n\ge2:\ \Sigma_\varepsilon^{(n)}\ \text{is connected}\}
\]
is an upper tail in \(\N\).  More precisely,
\[
\Sigma_\varepsilon^{(n)}\ \text{is connected}
\iff
\eta_n(a)<\varepsilon,
\]
so
\[
N_c(a,\varepsilon)=\min\{n\ge2:\ \eta_n(a)<\varepsilon\}
\]
satisfies
\[
n\ge N_c(a,\varepsilon)
\quad\Longrightarrow\quad
\Sigma_\varepsilon^{(n)}\ \text{is connected}.
\]
\end{corollary}

\begin{proof}
This is immediate from Proposition~\ref{prop:barrier}.
\end{proof}

Even without monotonicity of \(\eta_n(a)\), one gets a simple eventual connectedness
criterion from overlapping \(\varepsilon\)-discs around the eigenvalues.

\begin{lemma}[Eventual connectedness]\label{lem:eventual-connectedness}
If
\[
2\sqrt a\,\sin\!\left(\frac{\pi}{2(n+1)}\right)<\varepsilon,
\]
then \(\Sigma_\varepsilon^{(n)}\) is connected.

Consequently, for every fixed \(a\in(0,1)\) and \(\varepsilon>0\), the set
\[
\mathcal N_*(a,\varepsilon)
=
\left\{
n\ge2:\;
2\sqrt a\,\sin\!\left(\frac{\pi}{2(n+1)}\right)<\varepsilon
\right\}
\]
is nonempty, and with
\[
N_*(a,\varepsilon):=\min \mathcal N_*(a,\varepsilon),
\]
one has
\[
n\ge N_*(a,\varepsilon)
\quad\Longrightarrow\quad
\Sigma_\varepsilon^{(n)}\ \text{is connected}.
\]
\end{lemma}

\begin{proof}
Let
\[
\mu_k^{(n)}=2\sqrt a\,\cos\!\left(\frac{k\pi}{n+1}\right),
\qquad k=1,\dots,n,
\]
and let \(v_k\) be a corresponding unit eigenvector.  Then for every \(z\in\C\),
\[
s_n(z)\le \|(z\I_n-\A_n(a))v_k\|_2=|z-\mu_k^{(n)}|.
\]
Hence each open disc
\[
D_k:=\{z\in\C:\ |z-\mu_k^{(n)}|<\varepsilon\}
\]
satisfies \(D_k\subset \Sigma_\varepsilon^{(n)}\).

Now
\begin{align*}
\mu_k^{(n)}-\mu_{k+1}^{(n)}
&=
2\sqrt a
\left(
\cos\!\frac{k\pi}{n+1}
-
\cos\!\frac{(k+1)\pi}{n+1}
\right) \\
&=
4\sqrt a\,
\sin\!\left(\frac{(2k+1)\pi}{2(n+1)}\right)
\sin\!\left(\frac{\pi}{2(n+1)}\right) \\
&\le
4\sqrt a\,\sin\!\left(\frac{\pi}{2(n+1)}\right).
\end{align*}
Thus, if
\[
2\sqrt a\,\sin\!\left(\frac{\pi}{2(n+1)}\right)<\varepsilon,
\]
then
\[
\mu_k^{(n)}-\mu_{k+1}^{(n)}<2\varepsilon
\qquad(k=1,\dots,n-1),
\]
so
\[
D_k\cap D_{k+1}\neq\varnothing
\qquad(k=1,\dots,n-1).
\]
Therefore
\[
\mathcal D:=\bigcup_{k=1}^n D_k
\]
is connected.

By Lemma~\ref{lem:component-eigenvalue}, every connected component of
\(\Sigma_\varepsilon^{(n)}\) contains an eigenvalue of \(\A_n(a)\). Since the connected set
\(\mathcal D\subset \Sigma_\varepsilon^{(n)}\) contains all eigenvalues, there can be only
one connected component. Hence \(\Sigma_\varepsilon^{(n)}\) is connected.

Finally,
\[
\sin\!\left(\frac{\pi}{2(n+1)}\right)\to0
\qquad(n\to\infty),
\]
so \(\mathcal N_*(a,\varepsilon)\neq\varnothing\), and the last implication is immediate.
\end{proof}

\section{Singular vectors, currents, and continuation}
\label{sec:singular}

From now on we return to fixed \(n\ge2\) and \(0<a<1\), and use the shorthand
notation
\[
\A=\A_n(a),\qquad
\Bmat_{x,y}=(x+iy)\I_n-\A,\qquad
\Hmat_x(y)=\Bmat_{x,y}^*\Bmat_{x,y}.
\]
We also introduce the reversal matrix \(\PiM\in\R^{n\times n}\) by
\[
(\PiM)_{j,n+1-j}=1,\qquad (\PiM)_{jk}=0\ \text{otherwise}.
\]

\begin{lemma}[Canonical gauge]\label{lem:canonical}
One has
\[
\PiM\,\A\,\PiM=\A^{\mathsf T}.
\]
Suppose one of the variables \(x\) or \(y\) varies on an interval on which
\(\lambda_1(x,y)\) is simple. Then there exist real-analytic unit vectors
\(\bm u(x,y),\bm v(x,y)\in\C^n\) such that
\[
\Bmat_{x,y}\bm v=s(x,y)\bm u,
\qquad
\Bmat_{x,y}^*\bm u=s(x,y)\bm v,
\qquad
\bm u=\PiM\,\overline{\bm v}.
\]
\end{lemma}

\begin{proof}
The identity \(\PiM\A\PiM=\A^{\mathsf T}\) is immediate from the Toeplitz form of
\(\A\).

On an interval of simplicity, analytic perturbation theory yields a real-analytic unit
eigenvector \(\bm v\) of \(\Hmat_x(y)\) for the simple eigenvalue
\(\lambda_1(x,y)=s(x,y)^2\).  Since \(y>0\) and \(\sigma(\A)\subset\R\),
we have \(s(x,y)>0\).  Define
\[
\bm u=s(x,y)^{-1}\Bmat_{x,y}\bm v.
\]
Then \(\bm u\) is real-analytic and
\[
\Bmat_{x,y}\bm v=s(x,y)\bm u,
\qquad
\Bmat_{x,y}^*\bm u=s(x,y)\bm v.
\]

Because
\[
\PiM\Bmat_{x,y}=\Bmat_{x,y}^{\mathsf T}\PiM,
\]
the pair
\[
\PiM\,\overline{\bm u},
\qquad
\PiM\,\overline{\bm v}
\]
is another singular pair for the same simple singular value.  Hence there exists a
real-analytic unimodular scalar \(\alpha\) such that
\[
\PiM\,\overline{\bm u}=\alpha\,\bm v,
\qquad
\PiM\,\overline{\bm v}=\alpha\,\bm u.
\]
On an interval we may choose a real-analytic unimodular square root \(\beta\) of
\(\alpha\).  Replacing \((\bm u,\bm v)\) by \((\beta\bm u,\beta\bm v)\) yields
\[
\bm u=\PiM\,\overline{\bm v}.
\]
\end{proof}

\begin{proposition}[Current identity]\label{prop:current}
Fix \(x\in\R\), and let \(I\subset(0,\infty)\) be an interval on which
\(\lambda_1(x,y)\) is simple.  Let \(\bm v(y)\) be the canonical right singular vector
from Lemma~\ref{lem:canonical}, written as
\[
\bm v(y)=
\begin{bmatrix}
v_1(y)\\
\vdots\\
v_n(y)
\end{bmatrix},
\qquad
v_0(y)=v_{n+1}(y)=0.
\]
Define
\[
J_j(y)=\Im\!\bigl(\overline{v_j(y)}\,v_{j+1}(y)\bigr),
\qquad
j=1,\dots,n-1,
\]
and set \(J_0(y)=J_n(y)=0\).  Then, for \(j=1,\dots,n\),
\begin{equation}\label{eq:current}
J_j-aJ_{j-1}
=
y|v_j|^2+s(x,y)\,\Im\!\bigl(v_jv_{n+1-j}\bigr),
\end{equation}
and
\begin{equation}\label{eq:lambda-derivative}
\frac{d}{dy}\lambda_1(x,y)
=
2y-2(1-a)\sum_{j=1}^{n-1}J_j(y).
\end{equation}
\end{proposition}

\begin{proof}
Let
\[
\Bmat_y=(x+iy)\I_n-\A,\qquad s=s(x,y).
\]
Since
\[
\Bmat_y\bm v=s\bm u,
\qquad
\bm u=\PiM\,\overline{\bm v},
\]
the \(j\)-th component reads
\[
(x+iy)v_j-av_{j-1}-v_{j+1}
=
s\,\overline{v_{n+1-j}}.
\]
Multiplying by \(\overline{v_j}\) and taking imaginary parts gives
\[
y|v_j|^2-a\,\Im(\overline{v_j}v_{j-1})-\Im(\overline{v_j}v_{j+1})
=
-s\,\Im(v_jv_{n+1-j}),
\]
which is exactly \eqref{eq:current}.

For the derivative, the singular-value perturbation formula gives
\[
\frac{d}{dy}s(x,y)
=
\Re\!\bigl(\bm u^* i\bm v\bigr)
=
-\Im(\bm u^*\bm v).
\]
Since \(\bm u=\PiM\,\overline{\bm v}\),
\[
\bm u^*\bm v
=
\bm v^{\mathsf T}\PiM\bm v
=
\sum_{j=1}^n v_jv_{n+1-j}.
\]
Summing \eqref{eq:current} over \(j=1,\dots,n\), using \(J_0=J_n=0\), and
\(\|\bm v\|_2=1\), we obtain
\[
(1-a)\sum_{j=1}^{n-1}J_j
=
y+s\,\Im(\bm u^*\bm v).
\]
Therefore
\[
s\,\frac{d}{dy}s
=
y-(1-a)\sum_{j=1}^{n-1}J_j.
\]
Multiplying by \(2\) gives \eqref{eq:lambda-derivative}.
\end{proof}

\begin{lemma}[A real telescoping invariant]\label{lem:W}
Assume \(\lambda_1(x,y)\) is simple and let \(\bm v\) be in canonical gauge.
Define
\[
W_j=v_{j+1}v_{n+1-j}-a\,v_jv_{n-j},
\qquad
j=0,\dots,n,
\]
with \(v_0=v_{n+1}=0\).  Then
\[
W_0=0,
\qquad
W_j-W_{j-1}
=
s(x,y)\bigl(|v_j|^2-|v_{n+1-j}|^2\bigr)\in\R,
\]
so that
\[
W_j\in\R,\qquad j=0,\dots,n.
\]
\end{lemma}

\begin{proof}
Using the singular-vector equations at the indices \(j\) and \(n+1-j\),
\[
v_{j+1}+av_{j-1}=(x+iy)v_j-s\,\overline{v_{n+1-j}},
\]
\[
v_{n+2-j}+av_{n-j}=(x+iy)v_{n+1-j}-s\,\overline{v_j},
\]
we obtain
\begin{align*}
W_j-W_{j-1}
&=
v_{n+1-j}\bigl(v_{j+1}+av_{j-1}\bigr)
-
v_j\bigl(v_{n+2-j}+av_{n-j}\bigr) \\
&=
s\bigl(|v_j|^2-|v_{n+1-j}|^2\bigr)\in\R.
\end{align*}
Since \(W_0=0\), the claim follows.
\end{proof}

\begin{lemma}[Boundary entries on simple branches]\label{lem:boundary}
Assume \(\lambda_1(x,y)\) is simple and \(y>0\).  In canonical gauge,
\[
v_1\neq0,
\qquad
v_n\neq0.
\]
\end{lemma}

\begin{proof}
If \(n=2\), the claim is immediate because \(x+iy\neq0\).

Assume first that \(v_n=0\).  If also \(v_1=0\), then the equations at \(j=1\) and
\(j=n\) imply \(v_2=0\) and \(v_{n-1}=0\), and repeated propagation gives
\(\bm v=0\), impossible.  Hence \(v_1\neq0\).

From the equations at \(j=1\) and \(j=n-1\) we get
\[
v_2=(x+iy)v_1
\]
and
\[
(x+iy)v_{n-1}-av_{n-2}=s(x-iy)\overline{v_1}.
\]
Therefore
\[
W_{n-2}
=
v_{n-1}v_2-av_{n-2}v_1
=
s(x-iy)|v_1|^2.
\]
By Lemma~\ref{lem:W}, \(W_{n-2}\in\R\), which is impossible because \(y>0\) and
\(v_1\neq0\).  Thus \(v_n\neq0\).

The proof that \(v_1\neq0\) is analogous.  If \(v_1=0\), then
\[
v_2=-s\,\overline{v_n},
\qquad
av_{n-1}=(x+iy)v_n.
\]
Using the equation at \(j=2\), one obtains
\[
W_2
=
-\frac{s}{a}(x-iy)|v_n|^2,
\]
again contradicting Lemma~\ref{lem:W}.
\end{proof}

\begin{proposition}[Imaginary-axis seed]\label{prop:x0}
For every \(y>0\), the smallest eigenvalue \(\lambda_1(0,y)\) of \(\Hmat_0(y)\)
is simple.  In canonical gauge, the corresponding right singular vector can be written
as
\[
\bm v(0,y)
=
\beta\,
\begin{bmatrix}
w_1\\
iw_2\\
\vdots\\
i^{\,n-1}w_n
\end{bmatrix},
\qquad
w_j>0,
\]
for some unimodular scalar \(\beta\).  Consequently,
\[
\Im\!\left(\frac{v_{j+1}(0,y)}{v_j(0,y)}\right)>0,
\qquad
j=1,\dots,n-1,
\]
and
\[
\Im\!\left(\frac{\overline{v_{n+1-j}(0,y)}}{v_j(0,y)}\right)>0,
\qquad
j=1,\dots,n.
\]
\end{proposition}

\begin{proof}
Set
\[
\mathbf U=
\diag(1,i,i^2,\dots,i^{\,n-1}),
\qquad
\widehat{\Hmat}_y=\mathbf U^* \Hmat_0(y)\mathbf U.
\]
A direct computation shows that \(\widehat{\Hmat}_y\) is real symmetric, with diagonal
entries
\[
a^2+y^2,\quad 1+a^2+y^2,\dots,1+a^2+y^2,\quad 1+y^2,
\]
first off-diagonal entries
\[
-(1-a)y<0,
\]
and second off-diagonal entries
\[
-a<0.
\]
Hence \(\widehat{\Hmat}_y\) is a real symmetric irreducible \(Z\)-matrix.
Choose \(\Lambda>\max_j (\widehat{\Hmat}_y)_{jj}\).  Then
\(\Lambda \I_n-\widehat{\Hmat}_y\) is entrywise nonnegative and irreducible.  By
Perron--Frobenius, the smallest eigenvalue of \(\widehat{\Hmat}_y\) is simple and has a
strictly positive eigenvector
\[
\bm w=
\begin{bmatrix}
w_1\\
\vdots\\
w_n
\end{bmatrix},
\qquad
w_j>0.
\]
Therefore a right singular vector of \(\Bmat_{0,y}=iy\I_n-\A\) for \(s(0,y)\) can be
written as
\[
v_j=\beta\,i^{\,j-1}w_j,
\qquad
j=1,\dots,n,
\]
for some unimodular \(\beta\).

Insert this into
\[
iy\,v_j-av_{j-1}-v_{j+1}=s(0,y)\,\overline{v_{n+1-j}}.
\]
After division by \(\beta i^j\), we get
\[
w_{j+1}-yw_j-aw_{j-1}
=
-\gamma\,s(0,y)\,w_{n+1-j},
\qquad
\gamma=\frac{\overline\beta}{\beta}(-i)^n.
\]
At \(j=n\), this becomes
\[
yw_n+aw_{n-1}=\gamma\,s(0,y)\,w_1.
\]
The left-hand side is strictly positive, so \(\gamma=1\).

Now
\[
\frac{v_{j+1}(0,y)}{v_j(0,y)}
=
i\,\frac{w_{j+1}}{w_j},
\]
hence
\[
\Im\!\left(\frac{v_{j+1}(0,y)}{v_j(0,y)}\right)>0.
\]
Similarly,
\[
\frac{\overline{v_{n+1-j}(0,y)}}{v_j(0,y)}
=
i\,\frac{w_{n+1-j}}{w_j},
\]
which gives the second family of inequalities.
\end{proof}

\begin{lemma}[Angular identity]\label{lem:angular}
Assume \(\lambda_1(x,y)\) is simple and let \(\bm v\) be in canonical gauge.
For \(j=1,\dots,n-1\), define
\[
r_j=\frac{v_{j+1}}{v_j},
\qquad
\rho_j=r_{n-j},
\qquad
p_j=\frac{\overline{v_{n+1-j}}}{v_j},
\qquad
r_0=\infty,
\]
and
\[
\alpha_j=(x+iy)-\frac{a}{r_{j-1}},
\qquad
\beta_j=(x-iy)-\overline{\rho_{j-1}}.
\]
Then
\[
\Im\!\bigl(\alpha_j\overline{p_j}\bigr)
=
-\Im\!\bigl(\beta_j p_j\bigr).
\]
\end{lemma}

\begin{proof}
From the singular-vector equation at index \(j\),
\[
\left((x+iy)-\frac{a}{r_{j-1}}\right)v_j
=
v_{j+1}+s\,\overline{v_{n+1-j}},
\]
we obtain
\[
\alpha_j\overline{p_j}
=
\frac{v_{j+1}v_{n+1-j}+s|v_{n+1-j}|^2}{|v_j|^2},
\]
hence
\[
\Im\!\bigl(\alpha_j\overline{p_j}\bigr)
=
\frac{\Im(v_{j+1}v_{n+1-j})}{|v_j|^2}.
\]

Conjugating the singular-vector equation at index \(n+1-j\) yields
\[
(x-iy)\overline{v_{n+1-j}}
-
a\overline{v_{n-j}}
-
\overline{v_{n+2-j}}
=
sv_j.
\]
Therefore
\[
\beta_j p_j
=
s+a\,\frac{\overline{v_{n-j}}}{v_j},
\]
so
\[
\Im(\beta_j p_j)
=
-\frac{a\,\Im(v_jv_{n-j})}{|v_j|^2}.
\]
By Lemma~\ref{lem:W},
\[
\Im(v_{j+1}v_{n+1-j})=a\,\Im(v_jv_{n-j}),
\]
and the claim follows.
\end{proof}

\begin{proposition}[Phase continuation and branchwise monotonicity]
\label{prop:continuation}
Fix \(y>0\), and let \(I\subset\R\) be an open interval on which
\(\lambda_1(x,y)\) is simple.  Let
\[
\bm v(x)=
\begin{bmatrix}
v_1(x)\\
\vdots\\
v_n(x)
\end{bmatrix}
\]
be a real-analytic canonical right singular vector, and let \(J\) be a connected
component of
\[
\{x\in I:\ v_1(x)\cdots v_n(x)\neq0\}.
\]
Assume that for some \(x_0\in J\),
\[
\Im\!\left(\frac{v_{j+1}(x_0)}{v_j(x_0)}\right)>0,
\qquad j=1,\dots,n-1,
\]
and
\[
\Im\!\left(\frac{\overline{v_{n+1-j}(x_0)}}{v_j(x_0)}\right)>0,
\qquad j=1,\dots,n.
\]
Then the same inequalities hold for every \(x\in J\).  In particular,
\[
\partial_y\lambda_1(x,y)\ge0,
\qquad x\in J.
\]
\end{proposition}

\begin{proof}
\emph{Step 1: continuation of the ratios \(r_j=v_{j+1}/v_j\).}
Since \(v_j(x)\neq0\) on \(J\), we may write
\[
v_j(x)=t_j(x)e^{i\phi_j(x)},
\qquad
t_j(x)>0,
\]
with continuous arguments \(\phi_j\).  Set
\[
\delta_j(x)=\phi_{j+1}(x)-\phi_j(x),
\qquad j=1,\dots,n-1.
\]
Then
\[
\frac{v_{j+1}(x)}{v_j(x)}
=
\frac{t_{j+1}(x)}{t_j(x)}e^{i\delta_j(x)},
\]
so \(\Im(v_{j+1}/v_j)>0\) is equivalent to \(\delta_j\in(0,\pi)\).

Write
\[
c=(1+a)x-i(1-a)y=|c|e^{-i\vartheta},
\qquad
\vartheta\in(0,\pi).
\]
A direct computation gives
\[
\langle \bm v,\Hmat_x(y)\bm v\rangle
=
\sum_{j=1}^n \gamma_j t_j^2+\mathcal P_x(\delta_1,\dots,\delta_{n-1}),
\]
where \(\gamma_j\in\R\) and
\[
\mathcal P_x
=
-2|c|\sum_{j=1}^{n-1}t_jt_{j+1}\cos(\delta_j-\vartheta)
+
2a\sum_{j=1}^{n-2}t_jt_{j+2}\cos(\delta_j+\delta_{j+1}).
\]
For fixed amplitudes \(t_1,\dots,t_n\), the phase vector
\[
(\delta_1(x),\dots,\delta_{n-1}(x))
\]
minimizes \(\mathcal P_x\), because \(\bm v(x)\) is a ground state of
\(\Hmat_x(y)\).

Let
\[
E=
\{x\in J:\ \delta_j(x)\in(0,\pi)\ \text{for all }j=1,\dots,n-1\}.
\]
By assumption, \(x_0\in E\), so \(E\neq\varnothing\).  Clearly \(E\) is open in \(J\).
We show that \(E\) is closed in \(J\).

Take \(x_*\in\overline E\cap J\).  Then \(\delta_j(x_*)\in[0,\pi]\) for all \(j\).
Suppose that \(\delta_k(x_*)\in\{0,\pi\}\) for some \(k\).  Vary the tail phases by
\[
\phi_m\mapsto \phi_m+\varepsilon,\qquad m\ge k+1.
\]
This changes only \(\delta_k\), replacing it by \(\delta_k+\varepsilon\).  Hence the
derivative at \(\varepsilon=0\) of the phase functional equals
\begin{align*}
f_k'(0)
&=
2|c|\,t_kt_{k+1}\sin(\delta_k-\vartheta) \\
&\quad
-2a\,\mathbf 1_{\{k\ge2\}}\,t_{k-1}t_{k+1}\sin(\delta_{k-1}+\delta_k) \\
&\quad
-2a\,\mathbf 1_{\{k\le n-2\}}\,t_kt_{k+2}\sin(\delta_k+\delta_{k+1}).
\end{align*}
If \(\delta_k(x_*)=0\), then the first term is strictly negative and the other two are
nonpositive, so \(f_k'(0)<0\), contradicting minimality.  If \(\delta_k(x_*)=\pi\),
then the first term is strictly positive and the other two are nonnegative, so
\(f_k'(0)>0\), again contradicting minimality.  Therefore no \(\delta_k\) can hit
\(0\) or \(\pi\), and thus \(x_*\in E\).  Hence \(E\) is closed in \(J\).

Since \(J\) is connected and \(E\) is nonempty, open, and closed in \(J\), we conclude
that \(E=J\).  Thus
\[
\Im\!\left(\frac{v_{j+1}(x)}{v_j(x)}\right)>0,
\qquad
j=1,\dots,n-1,\ \ x\in J.
\]
In particular,
\[
\Im \rho_j(x)>0,
\qquad
j=1,\dots,n-1,\ \ x\in J.
\]

\medskip
\noindent
\emph{Step 2: continuation of the reflected phases
\(p_j=\overline{v_{n+1-j}}/v_j\).}
For \(j=1,\dots,n-1\), suppose that \(\Im p_j\) vanishes somewhere on \(J\). By
continuity there exists \(x_*\in J\) with \(p_j(x_*)\in\R\). Since all components of
\(\bm v\) are nonzero on \(J\), we have \(p_j(x_*)\neq0\).

By Step 1,
\[
\Im\alpha_j(x_*)
=
y+\Im\!\left(-\frac{a}{r_{j-1}(x_*)}\right)>0.
\]
Lemma~\ref{lem:angular} therefore gives
\[
\Im\beta_j(x_*)=-\Im\alpha_j(x_*)<0.
\]
Conjugating the singular-vector equation at index \(n+1-j\) shows that
\[
\overline{\rho_j}
=
\frac{a}{\beta_j-s/p_j}.
\]
Since \(s/p_j\in\R\) at \(x_*\), the denominator has negative imaginary part, and
hence \(\Im\overline{\rho_j}(x_*)>0\), i.e.,
\[
\Im\rho_j(x_*)<0,
\]
contradicting Step 1.  Thus
\[
\Im p_j(x)>0,
\qquad
j=1,\dots,n-1,\ \ x\in J.
\]
For \(j=n\) we use
\[
p_n(x)=\frac{1}{\overline{p_1(x)}},
\]
which shows \(\Im p_n(x)>0\) as well.

\medskip
\noindent
\emph{Step 3: monotonicity in \(y\).}
Since
\[
p_j
=
\frac{\overline{v_{n+1-j}}}{v_j}
=
\frac{\overline{v_jv_{n+1-j}}}{|v_j|^2},
\]
we have
\[
\Im p_j
=
-\frac{\Im(v_jv_{n+1-j})}{|v_j|^2}.
\]
Therefore \(\Im p_j>0\) implies
\[
\Im(v_jv_{n+1-j})<0.
\]
Using \eqref{eq:current}, we obtain
\[
J_j-aJ_{j-1}\le y|v_j|^2,
\qquad
j=1,\dots,n.
\]
Inductively,
\[
J_j
\le
y\sum_{k=1}^j a^{\,j-k}|v_k|^2,
\qquad
j=1,\dots,n-1.
\]
Summing this inequality yields
\[
(1-a)\sum_{j=1}^{n-1}J_j
\le
y\sum_{k=1}^{n-1}(1-a^{\,n-k})|v_k|^2
\le y.
\]
Now \eqref{eq:lambda-derivative} gives
\[
\partial_y\lambda_1(x,y)\ge0,
\qquad x\in J.
\]
\end{proof}

\section{One-sided positivity and edge Schur complements}
\label{sec:schur}

Fix \(x\in\R\) and \(y>0\), and write
\[
z=x+iy,
\qquad
\lambda=\lambda_1(x,y),
\qquad
c=\overline z+az=(1+a)x-i(1-a)y.
\]
Set
\[
d_-=a^2+|z|^2-\lambda,
\qquad
d=1+a^2+|z|^2-\lambda,
\qquad
d_+=1+|z|^2-\lambda.
\]
Then
\[
\Hmat_x(y)-\lambda\I_n
=
\begin{bmatrix}
d_- & -c & a \\
-\overline c & d & -c & a \\
a & -\overline c & d & -c & a \\
& \ddots & \ddots & \ddots & \ddots & \ddots \\
&& a & -\overline c & d & -c \\
&&& a & -\overline c & d_+
\end{bmatrix},
\]
with all omitted entries equal to zero.  For an index set
\(I\subset\{1,\dots,n\}\), we write \(\mathbf M[I]\) for the principal submatrix of
\(\mathbf M\) indexed by \(I\).

For \(j=2,\dots,n-1\), define
\[
\mathbf L_j
=
\bigl(\Hmat_x(y)-\lambda\I_n\bigr)[\{1,\dots,j-1\}],
\qquad
\mathbf R_j
=
\bigl(\Hmat_x(y)-\lambda\I_n\bigr)[\{j+1,\dots,n\}].
\]

\begin{proposition}[One-sided positivity]\label{prop:onesided}
For every \(j=2,\dots,n-1\),
\[
\mathbf L_j\succ0,
\qquad
\mathbf R_j\succ0.
\]
In particular,
\[
d_->0,
\qquad
d_+>0.
\]
\end{proposition}

\begin{proof}
Since \(\Hmat_x(y)-\lambda\I_n\succeq0\), every principal submatrix is positive
semidefinite.  It therefore suffices to prove that the kernels of \(\mathbf L_j\) and
\(\mathbf R_j\) are trivial.

We prove this for \(\mathbf R_j\); the proof for \(\mathbf L_j\) is symmetric.
Let \(\bm w\in\ker\mathbf R_j\), and extend it by zeros to
\[
\widehat{\bm w}=
\begin{bmatrix}
0\\
\vdots\\
0\\
w_1\\
\vdots\\
w_{n-j}
\end{bmatrix}\in\C^n.
\]
Then rows \(j+1,\dots,n\) of
\((\Hmat_x(y)-\lambda\I_n)\widehat{\bm w}\) vanish.  Row \(j-1\) gives
\[
aw_1=0,
\]
hence \(w_1=0\).  If \(n-j\ge2\), then row \(j\) gives
\[
aw_2=0,
\]
hence \(w_2=0\).  The first row of \(\mathbf R_j\bm w=0\) now yields
\[
aw_3=0.
\]
Proceeding inductively, all components vanish.  Thus
\(\ker\mathbf R_j=\{0\}\), so \(\mathbf R_j\succ0\).

The proof for \(\mathbf L_j\) is identical.  Finally,
\[
\mathbf L_2=[d_-],
\qquad
\mathbf R_{n-1}=[d_+],
\]
hence \(d_->0\) and \(d_+>0\).
\end{proof}

\begin{definition}\label{def:edge}
Define
\[
\mathbf T=
\begin{bmatrix}
d & -c\\
-\overline c & d
\end{bmatrix},
\qquad
\mathbf E=
\begin{bmatrix}
a & 0\\
-c & a
\end{bmatrix}.
\]
For \(j=3,\dots,n-1\), let \(\mathbf S_j^{\mathrm L}\) be the Schur complement of
\(\mathbf L_j\) onto its last two coordinates \((j-2,j-1)\).  For
\(j=2,\dots,n-2\), let \(\mathbf S_j^{\mathrm R}\) be the Schur complement of
\(\mathbf R_j\) onto its first two coordinates \((j+1,j+2)\).
\end{definition}

\begin{proposition}[Riccati recursions]\label{prop:schur-rec}
Whenever the indices are meaningful,
\[
\mathbf S_3^{\mathrm L}
=
\begin{bmatrix}
d_- & -c\\
-\overline c & d
\end{bmatrix},
\qquad
\mathbf S_4^{\mathrm L}
=
\begin{bmatrix}
d-\dfrac{|c|^2}{d_-} & -c+\dfrac{a\overline c}{d_-}\\[3mm]
-\overline c+\dfrac{ac}{d_-} & d-\dfrac{a^2}{d_-}
\end{bmatrix},
\]
and
\[
\mathbf S_j^{\mathrm L}
=
\mathbf T-\mathbf E^*(\mathbf S_{j-2}^{\mathrm L})^{-1}\mathbf E,
\qquad
j\ge5.
\]
Similarly,
\[
\mathbf S_{n-2}^{\mathrm R}
=
\begin{bmatrix}
d & -c\\
-\overline c & d_+
\end{bmatrix},
\qquad
\mathbf S_{n-3}^{\mathrm R}
=
\begin{bmatrix}
d-\dfrac{a^2}{d_+} & -c+\dfrac{a\overline c}{d_+}\\[3mm]
-\overline c+\dfrac{ac}{d_+} & d-\dfrac{|c|^2}{d_+}
\end{bmatrix},
\]
and
\[
\mathbf S_j^{\mathrm R}
=
\mathbf T-\mathbf E(\mathbf S_{j+2}^{\mathrm R})^{-1}\mathbf E^*,
\qquad
j\le n-4.
\]
\end{proposition}

\begin{proof}
These formulas follow from direct block Schur-complement calculations for the
five-diagonal matrix \(\Hmat_x(y)-\lambda\I_n\).
\end{proof}

\begin{lemma}[Propagation of the upper half-plane]\label{lem:propagate}
Let
\[
\mathbf S=
\begin{bmatrix}
\alpha & \beta\\
\overline\beta & \delta
\end{bmatrix}\succ0,
\qquad
\Delta=\alpha\delta-|\beta|^2>0.
\]
If
\[
\mathbf T_{\mathrm L}
=
\mathbf T-\mathbf E^*\mathbf S^{-1}\mathbf E
=
\begin{bmatrix}
\alpha' & \beta'\\
\overline{\beta'} & \delta'
\end{bmatrix},
\]
then
\[
\beta'
=
\frac{a^2\beta+a\alpha\,\overline c-c\,\Delta}{\Delta},
\]
and therefore
\[
\Im\beta'
=
\frac{a^2\Im\beta+a\alpha(1-a)y+\Delta(1-a)y}{\Delta}.
\]
If
\[
\mathbf T_{\mathrm R}
=
\mathbf T-\mathbf E\mathbf S^{-1}\mathbf E^*
=
\begin{bmatrix}
\widetilde\alpha & \widetilde\beta\\
\overline{\widetilde\beta} & \widetilde\delta
\end{bmatrix},
\]
then
\[
\widetilde\beta
=
\frac{a^2\beta+a\delta\,\overline c-c\,\Delta}{\Delta},
\]
and therefore
\[
\Im\widetilde\beta
=
\frac{a^2\Im\beta+a\delta(1-a)y+\Delta(1-a)y}{\Delta}.
\]
In particular, if \(\Im\beta>0\), then
\[
\Im\beta'>0,
\qquad
\Im\widetilde\beta>0.
\]
\end{lemma}

\begin{proof}
Using
\[
\mathbf S^{-1}
=
\frac{1}{\Delta}
\begin{bmatrix}
\delta & -\beta\\
-\overline\beta & \alpha
\end{bmatrix},
\]
one computes the off-diagonal entries of \(\mathbf T_{\mathrm L}\) and
\(\mathbf T_{\mathrm R}\) directly.
\end{proof}

\begin{corollary}[Upper-half-plane property]\label{cor:upper}
For every \(j=3,\dots,n-1\),
\[
\mathbf S_j^{\mathrm L}
=
\begin{bmatrix}
\alpha_j^{\mathrm L} & \beta_j^{\mathrm L}\\
\overline{\beta_j^{\mathrm L}} & \delta_j^{\mathrm L}
\end{bmatrix}
\quad\Longrightarrow\quad
\Im\beta_j^{\mathrm L}>0.
\]
For every \(j=2,\dots,n-2\),
\[
\mathbf S_j^{\mathrm R}
=
\begin{bmatrix}
\alpha_j^{\mathrm R} & \beta_j^{\mathrm R}\\
\overline{\beta_j^{\mathrm R}} & \delta_j^{\mathrm R}
\end{bmatrix}
\quad\Longrightarrow\quad
\Im\beta_j^{\mathrm R}>0.
\]
\end{corollary}

\begin{proof}
The base cases are
\[
\beta_3^{\mathrm L}=-c,
\qquad
\beta_4^{\mathrm L}=-c+\frac{a\overline c}{d_-},
\qquad
\beta_{n-2}^{\mathrm R}=-c,
\qquad
\beta_{n-3}^{\mathrm R}=-c+\frac{a\overline c}{d_+},
\]
all of which have positive imaginary part because \(d_\pm>0\).  The recursions in
Proposition~\ref{prop:schur-rec} and Lemma~\ref{lem:propagate} complete the proof.
\end{proof}

\section{A no-zero theorem for ground states}
\label{sec:nozero}

\begin{proposition}[The second boundary layer never vanishes]
\label{prop:second-layer}
Assume \(n\ge3\), and let
\[
\bm v=
\begin{bmatrix}
v_1\\
\vdots\\
v_n
\end{bmatrix}
\in
\ker\bigl(\Hmat_x(y)-\lambda\I_n\bigr)\setminus\{0\},
\qquad
\lambda=\lambda_1(x,y).
\]
Then
\[
v_2\neq0,
\qquad
v_{n-1}\neq0.
\]
\end{proposition}

\begin{proof}
We prove \(v_2\neq0\); the other statement follows by reversing the coordinates.

Assume \(v_2=0\).

If \(n=3\), then \(v_1\neq0\); otherwise \(\mathbf R_2[v_3]^{\mathsf T}=0\) would
contradict Proposition~\ref{prop:onesided}.  Row \(3\) gives
\[
d_+v_3+av_1=0,
\qquad\text{so}\qquad
v_3=-\frac{a}{d_+}v_1.
\]
Row \(2\) then gives
\[
-\overline c\,v_1-c\,v_3=0,
\]
hence
\[
\overline c=\frac{a}{d_+}c,
\]
which is impossible because \(a/d_+>0\), while
\(\Im(\overline c)=+(1-a)y\) and \(\Im(c)=-(1-a)y\).

Assume now that \(n\ge4\).  If \(v_1=0\), then
\[
\mathbf R_2
\begin{bmatrix}
v_3\\
\vdots\\
v_n
\end{bmatrix}
=
\bm0,
\]
contradicting Proposition~\ref{prop:onesided}.  Hence \(v_1\neq0\).

Let
\[
\mathbf G_2=\mathbf R_2^{-1},
\qquad
g=\langle \bm e_1,\mathbf G_2\bm e_1\rangle>0,
\qquad
h=\langle \bm e_2,\mathbf G_2\bm e_1\rangle,
\]
where \(\bm e_1,\bm e_2\) are the first two basis vectors of \(\C^{n-2}\).
From
\[
\mathbf R_2
\begin{bmatrix}
v_3\\
\vdots\\
v_n
\end{bmatrix}
+
av_1\bm e_1
=
\bm0
\]
we obtain
\[
v_3=-ag\,v_1,
\qquad
v_4=-ah\,v_1.
\]
Row \(1\) of
\((\Hmat_x(y)-\lambda\I_n)\bm v=\bm0\) gives
\[
d_-v_1+av_3=0,
\]
so
\[
d_-=a^2g.
\]
Row \(2\) gives
\[
-\overline c\,v_1-c\,v_3+av_4=0,
\]
that is,
\[
-\overline c+ag\,c-a^2h=0.
\]
Equivalently,
\[
-c+a\frac{h}{g}
=
-\frac{1}{ag}\,\overline c
=
-\frac{a}{d_-}\,\overline c.
\]
Hence the left-hand side has negative imaginary part.

On the other hand, write
\[
\mathbf S_2^{\mathrm R}
=
\begin{bmatrix}
\alpha_2^{\mathrm R} & \beta_2^{\mathrm R}\\
\overline{\beta_2^{\mathrm R}} & \delta_2^{\mathrm R}
\end{bmatrix}.
\]
By the block inversion formula,
\[
\frac{h}{g}
=
-\frac{\overline{\beta_2^{\mathrm R}}}{\delta_2^{\mathrm R}},
\]
and therefore
\[
-c+a\frac{h}{g}
=
-c-a\,\frac{\overline{\beta_2^{\mathrm R}}}{\delta_2^{\mathrm R}}.
\]
By Corollary~\ref{cor:upper}, \(\Im\beta_2^{\mathrm R}>0\), so this number has positive
imaginary part, a contradiction.  Thus \(v_2\neq0\).
\end{proof}

\begin{proposition}[Defect resonance for an interior zero]\label{prop:defect}
Assume \(n\ge5\), and let
\[
\bm v=
\begin{bmatrix}
v_1\\
\vdots\\
v_n
\end{bmatrix}
\in
\ker\bigl(\Hmat_x(y)-\lambda\I_n\bigr)\setminus\{0\},
\qquad
\lambda=\lambda_1(x,y).
\]
Fix \(j\in\{3,\dots,n-2\}\) and assume \(v_j=0\).
Set
\[
\mathbf G_j^{\mathrm L}=\mathbf L_j^{-1},
\qquad
\mathbf G_j^{\mathrm R}=\mathbf R_j^{-1},
\]
and define
\[
g_j^{\mathrm L}
=
\langle \bm e_{j-1},\mathbf G_j^{\mathrm L}\bm e_{j-1}\rangle>0,
\qquad
h_j^{\mathrm L}
=
\langle \bm e_{j-2},\mathbf G_j^{\mathrm L}\bm e_{j-1}\rangle,
\]
where \(\bm e_{j-2},\bm e_{j-1}\) are the last two basis vectors of \(\C^{j-1}\), and
\[
g_j^{\mathrm R}
=
\langle \bm e_1,\mathbf G_j^{\mathrm R}\bm e_1\rangle>0,
\qquad
h_j^{\mathrm R}
=
\langle \bm e_2,\mathbf G_j^{\mathrm R}\bm e_1\rangle,
\]
where \(\bm e_1,\bm e_2\) are the first two basis vectors of \(\C^{n-j}\).
Then:
\begin{enumerate}
\item
\[
v_{j-1}=-ag_j^{\mathrm L}v_{j+1},
\qquad
v_{j+1}=-ag_j^{\mathrm R}v_{j-1}.
\]
In particular,
\[
v_{j-1}\neq0,
\qquad
v_{j+1}\neq0,
\qquad
a^2g_j^{\mathrm L}g_j^{\mathrm R}=1.
\]

\item
If
\[
m_j^{\mathrm L}=a\frac{h_j^{\mathrm L}}{g_j^{\mathrm L}}-\overline c,
\qquad
m_j^{\mathrm R}=-c+a\frac{h_j^{\mathrm R}}{g_j^{\mathrm R}},
\]
then
\[
m_j^{\mathrm L}v_{j-1}+m_j^{\mathrm R}v_{j+1}=0,
\]
hence
\[
m_j^{\mathrm R}=ag_j^{\mathrm L}m_j^{\mathrm L}.
\]
\end{enumerate}
\end{proposition}

\begin{proof}
Split
\[
\bm v=
\begin{bmatrix}
v_1\\
\vdots\\
v_{j-1}\\
0\\
v_{j+1}\\
\vdots\\
v_n
\end{bmatrix}
=
\begin{bmatrix}
\bm u\\
0\\
\bm w
\end{bmatrix}.
\]
The equations on the left and right of the defect read
\[
\mathbf L_j\bm u+av_{j+1}\bm e_{j-1}=\bm0,
\qquad
\mathbf R_j\bm w+av_{j-1}\bm e_1=\bm0.
\]
Therefore
\[
\bm u=-av_{j+1}\mathbf G_j^{\mathrm L}\bm e_{j-1},
\qquad
\bm w=-av_{j-1}\mathbf G_j^{\mathrm R}\bm e_1,
\]
which gives
\[
v_{j-1}=-ag_j^{\mathrm L}v_{j+1},
\qquad
v_{j+1}=-ag_j^{\mathrm R}v_{j-1}.
\]
If \(v_{j+1}=0\), then \(\bm u=\bm0\), hence \(v_{j-1}=0\), and then \(\bm w=\bm0\),
contrary to \(\bm v\neq0\).  Thus \(v_{j\pm1}\neq0\), and
\(a^2g_j^{\mathrm L}g_j^{\mathrm R}=1\) follows immediately.

Moreover,
\[
v_{j-2}=\frac{h_j^{\mathrm L}}{g_j^{\mathrm L}}v_{j-1},
\qquad
v_{j+2}=\frac{h_j^{\mathrm R}}{g_j^{\mathrm R}}v_{j+1}.
\]
The \(j\)-th row of
\((\Hmat_x(y)-\lambda\I_n)\bm v=\bm0\) is
\[
av_{j-2}-\overline c\,v_{j-1}-c\,v_{j+1}+av_{j+2}=0,
\]
which is equivalent to
\[
m_j^{\mathrm L}v_{j-1}+m_j^{\mathrm R}v_{j+1}=0.
\]
Using \(v_{j-1}=-ag_j^{\mathrm L}v_{j+1}\) and \(v_{j+1}\neq0\) yields
\[
m_j^{\mathrm R}=ag_j^{\mathrm L}m_j^{\mathrm L}.
\]
\end{proof}

\begin{proposition}[Sign separation for the defect Weyl quotients]
\label{prop:sign}
For \(j=3,\dots,n-2\), write
\[
\mathbf S_j^{\mathrm L}
=
\begin{bmatrix}
\alpha_j^{\mathrm L} & \beta_j^{\mathrm L}\\
\overline{\beta_j^{\mathrm L}} & \delta_j^{\mathrm L}
\end{bmatrix},
\qquad
\mathbf S_j^{\mathrm R}
=
\begin{bmatrix}
\alpha_j^{\mathrm R} & \beta_j^{\mathrm R}\\
\overline{\beta_j^{\mathrm R}} & \delta_j^{\mathrm R}
\end{bmatrix}.
\]
Then
\[
\frac{h_j^{\mathrm L}}{g_j^{\mathrm L}}
=
-\frac{\beta_j^{\mathrm L}}{\alpha_j^{\mathrm L}},
\qquad
\frac{h_j^{\mathrm R}}{g_j^{\mathrm R}}
=
-\frac{\overline{\beta_j^{\mathrm R}}}{\delta_j^{\mathrm R}},
\]
and therefore
\[
m_j^{\mathrm L}
=
-\overline c-a\,\frac{\beta_j^{\mathrm L}}{\alpha_j^{\mathrm L}},
\qquad
m_j^{\mathrm R}
=
-c-a\,\frac{\overline{\beta_j^{\mathrm R}}}{\delta_j^{\mathrm R}}.
\]
In particular,
\[
\Im m_j^{\mathrm L}<0<\Im m_j^{\mathrm R}.
\]
\end{proposition}

\begin{proof}
Since \(\mathbf S_j^{\mathrm L}\) is the Schur complement of \(\mathbf L_j\) onto the last
two coordinates, the lower-right \(2\times2\) block of \(\mathbf L_j^{-1}\) is
\((\mathbf S_j^{\mathrm L})^{-1}\).  Hence
\[
g_j^{\mathrm L}=\bigl((\mathbf S_j^{\mathrm L})^{-1}\bigr)_{22},
\qquad
h_j^{\mathrm L}=\bigl((\mathbf S_j^{\mathrm L})^{-1}\bigr)_{12}.
\]
Since
\[
(\mathbf S_j^{\mathrm L})^{-1}
=
\frac{1}{\alpha_j^{\mathrm L}\delta_j^{\mathrm L}-|\beta_j^{\mathrm L}|^2}
\begin{bmatrix}
\delta_j^{\mathrm L} & -\beta_j^{\mathrm L}\\
-\overline{\beta_j^{\mathrm L}} & \alpha_j^{\mathrm L}
\end{bmatrix},
\]
the first ratio formula follows.  The proof for \(\mathbf S_j^{\mathrm R}\) is analogous.

By Corollary~\ref{cor:upper},
\[
\Im\beta_j^{\mathrm L}>0,
\qquad
\Im\beta_j^{\mathrm R}>0.
\]
Since \(\alpha_j^{\mathrm L}>0\) and \(\delta_j^{\mathrm R}>0\), we get
\[
\Im m_j^{\mathrm L}
=
-(1-a)y-a\,\frac{\Im\beta_j^{\mathrm L}}{\alpha_j^{\mathrm L}}<0
\]
and
\[
\Im m_j^{\mathrm R}
=
(1-a)y+a\,\frac{\Im\beta_j^{\mathrm R}}{\delta_j^{\mathrm R}}>0.
\]
\end{proof}

\begin{theorem}[No interior zeros for a ground state]\label{thm:nozero}
Assume \(n\ge3\), and let
\[
\bm v\in
\ker\bigl(\Hmat_x(y)-\lambda_1(x,y)\I_n\bigr)\setminus\{0\}.
\]
Then
\[
v_j\neq0,
\qquad
j=2,\dots,n-1.
\]
\end{theorem}

\begin{proof}
By Proposition~\ref{prop:second-layer},
\[
v_2\neq0,
\qquad
v_{n-1}\neq0.
\]
Thus only the indices \(j=3,\dots,n-2\) remain.  If such a \(j\) satisfied \(v_j=0\),
then Proposition~\ref{prop:defect} would imply
\[
m_j^{\mathrm R}=ag_j^{\mathrm L}m_j^{\mathrm L}
\]
with \(ag_j^{\mathrm L}>0\).  Hence \(m_j^{\mathrm R}\) would be a positive real multiple of
\(m_j^{\mathrm L}\), contradicting Proposition~\ref{prop:sign}.
\end{proof}

\begin{corollary}[Global simplicity of the lowest branch]\label{cor:simple}
Assume \(n\ge3\).  Then for every \(x\in\R\) and \(y>0\), the smallest eigenvalue
\[
\lambda_1(x,y)=\lambda_{\min}\!\bigl(\Hmat_x(y)\bigr)
\]
is simple.
\end{corollary}

\begin{proof}
If the eigenspace had dimension at least \(2\), then the linear map
\[
\ker\bigl(\Hmat_x(y)-\lambda_1(x,y)\I_n\bigr)\longrightarrow\C,
\qquad
\bm v\longmapsto v_2,
\]
would have a nontrivial kernel.  This would contradict
Theorem~\ref{thm:nozero}.
\end{proof}

\section{Proof of vertical monotonicity and simple connectedness}
\label{sec:proof-vertical}

\begin{proposition}[The case \(n=2\)]\label{prop:n2}
If \(n=2\), then for all \(x\in\R\) and \(y\ge0\),
\[
\partial_y\lambda_1(x,y)\ge0.
\]
\end{proposition}

\begin{proof}
For \(n=2\),
\[
\Bmat_{x,y}
=
\begin{bmatrix}
x+iy & -1\\
-a & x+iy
\end{bmatrix},
\]
and hence
\[
\Hmat_x(y)
=
\begin{bmatrix}
a^2+x^2+y^2 & -(a+1)x+i(1-a)y\\
-(a+1)x-i(1-a)y & 1+x^2+y^2
\end{bmatrix}.
\]
Its smaller eigenvalue is
\[
\lambda_1(x,y)
=
x^2+y^2+\frac{1+a^2}{2}
-\frac12\sqrt{(1-a^2)^2+4(a+1)^2x^2+4(1-a)^2y^2}.
\]
Differentiating gives
\[
\partial_y\lambda_1(x,y)
=
2y\left(
1-\frac{(1-a)^2}
{\sqrt{(1-a^2)^2+4(a+1)^2x^2+4(1-a)^2y^2}}
\right)\ge0,
\]
because
\[
\sqrt{(1-a^2)^2+4(a+1)^2x^2+4(1-a)^2y^2}\ge 1-a^2\ge (1-a)^2.
\]
\end{proof}

\begin{proof}[Proof of Theorem~\ref{thm:vertical}]
If \(n=2\), the claim is Proposition~\ref{prop:n2}.  Assume henceforth that
\(n\ge3\).

Fix \(y_*>0\).  By Corollary~\ref{cor:simple}, the lowest eigenvalue
\[
x\longmapsto \lambda_1(x,y_*)
\]
is simple for every \(x\in\R\).  Since \(x\mapsto \Hmat_x(y_*)\) is a real-analytic
family, Lemma~\ref{lem:canonical} yields a real-analytic canonical right singular
vector
\[
\bm v(x)=
\begin{bmatrix}
v_1(x)\\
\vdots\\
v_n(x)
\end{bmatrix},
\qquad
x\in\R.
\]

By Theorem~\ref{thm:nozero},
\[
v_j(x)\neq0,
\qquad
j=2,\dots,n-1,\ \ x\in\R.
\]
By Lemma~\ref{lem:boundary}, also
\[
v_1(x)\neq0,
\qquad
v_n(x)\neq0,
\qquad
x\in\R.
\]
Thus every component of \(\bm v(x)\) is nonzero for every \(x\in\R\).

At \(x=0\), Proposition~\ref{prop:x0} gives
\[
\Im\!\left(\frac{v_{j+1}(0)}{v_j(0)}\right)>0,
\qquad
j=1,\dots,n-1,
\]
and
\[
\Im\!\left(\frac{\overline{v_{n+1-j}(0)}}{v_j(0)}\right)>0,
\qquad
j=1,\dots,n.
\]
Hence Proposition~\ref{prop:continuation}, applied with \(J=\R\), yields
\[
\partial_y\lambda_1(x,y_*)\ge0,
\qquad
x\in\R.
\]
Since \(y_*>0\) was arbitrary, we conclude that
\[
\partial_y\lambda_1(x,y)\ge0,
\qquad
x\in\R,\ y>0.
\]

Finally, Lemma~\ref{lem:conj} implies that \(\lambda_1(x,y)\) is even in \(y\).
Therefore \(y\mapsto \lambda_1(x,y)\) is nondecreasing on \([0,\infty)\), and since
\(s(x,y)=\sqrt{\lambda_1(x,y)}\ge0\), the same is true of \(y\mapsto s(x,y)\).
\end{proof}

\begin{proof}[Proof of Theorem~\ref{thm:main}]
Combine Theorem~\ref{thm:vertical} with
Proposition~\ref{prop:topological-reduction}.
\end{proof}

\begin{remark}
With Theorem~\ref{thm:vertical} now established, all the conditional results in
Section~\ref{sec:topology} are unconditional.  In particular,
Proposition~\ref{prop:barrier} and Corollary~\ref{cor:threshold} are proved.
\end{remark}

\appendix

\section{Real-axis factorization and boundary algebra}
\label{app:factor}

Let \(\PiM_n\in\R^{n\times n}\) be the reversal matrix,
\[
(\PiM_n)_{j,n+1-j}=1,\qquad (\PiM_n)_{jk}=0\ \text{otherwise}.
\]

\begin{proposition}[Reversal symmetry and factorization]\label{prop:factorization}
For every \(n\ge2\) and every real \(x\), one has
\[
\PiM_n\,\A_n(a)\,\PiM_n=\A_n(a)^{\mathsf T}.
\]
Hence the matrix
\[
\mathbf G_n(x):=(x\I_n-\A_n(a))\PiM_n
\]
is real symmetric. Moreover, for every \(\varepsilon\ge0\),
\begin{equation}\label{eq:factorization}
\det\!\bigl((x\I_n-\A_n(a))^{\mathsf T}(x\I_n-\A_n(a))-\varepsilon^2\I_n\bigr)
=
\det\!\bigl(x\I_n-(\A_n(a)+\varepsilon\PiM_n)\bigr)\,
\det\!\bigl(x\I_n-(\A_n(a)-\varepsilon\PiM_n)\bigr).
\end{equation}
\end{proposition}

\begin{proof}
The identity
\[
\PiM_n\,\A_n(a)\,\PiM_n=\A_n(a)^{\mathsf T}
\]
is immediate from the tridiagonal Toeplitz structure.

Using \(\PiM_n^2=\I_n\), we compute
\[
\mathbf G_n(x)^{\mathsf T}
=
\PiM_n(x\I_n-\A_n(a))^{\mathsf T}
=
\PiM_n(x\I_n-\A_n(a)^{\mathsf T})
=
(x\I_n-\A_n(a))\PiM_n
=
\mathbf G_n(x),
\]
so \(\mathbf G_n(x)\) is symmetric.

Next,
\[
\PiM_n\,(x\I_n-\A_n(a))^{\mathsf T}(x\I_n-\A_n(a))\,\PiM_n
=
(x\I_n-\A_n(a))(x\I_n-\A_n(a))^{\mathsf T}
=
\mathbf G_n(x)^2.
\]
Since \(\PiM_n\) is orthogonal and involutory,
\[
\det\!\bigl((x\I_n-\A_n(a))^{\mathsf T}(x\I_n-\A_n(a))-\varepsilon^2\I_n\bigr)
=
\det\!\bigl(\mathbf G_n(x)^2-\varepsilon^2\I_n\bigr).
\]
Because \(\mathbf G_n(x)\) is symmetric, the factors commute, and therefore
\[
\det\!\bigl(\mathbf G_n(x)^2-\varepsilon^2\I_n\bigr)
=
\det\!\bigl(\mathbf G_n(x)-\varepsilon\I_n\bigr)\,
\det\!\bigl(\mathbf G_n(x)+\varepsilon\I_n\bigr).
\]
Finally,
\[
\mathbf G_n(x)-\varepsilon\I_n
=
(x\I_n-\A_n(a))\PiM_n-\varepsilon\I_n
=
\bigl(x\I_n-\A_n(a)-\varepsilon\PiM_n\bigr)\PiM_n,
\]
and similarly
\[
\mathbf G_n(x)+\varepsilon\I_n
=
\bigl(x\I_n-\A_n(a)+\varepsilon\PiM_n\bigr)\PiM_n.
\]
Taking determinants and using \(\det(\PiM_n)^2=1\) yields
\eqref{eq:factorization}.
\end{proof}

\begin{corollary}[Real-axis boundary points come from two explicit spectral problems]
\label{cor:boundary}
Fix \(n\ge2\) and \(\varepsilon>0\).  If
\[
x\in\partial I_\varepsilon^{(n)},
\]
then
\[
x\in \sigma\!\bigl(\A_n(a)+\varepsilon\PiM_n\bigr)
\cup
\sigma\!\bigl(\A_n(a)-\varepsilon\PiM_n\bigr).
\]
Equivalently, every real boundary point of the \(\varepsilon\)-pseudospectrum is a real
eigenvalue of one of the two matrices \(\A_n(a)\pm \varepsilon\PiM_n\).
\end{corollary}

\begin{proof}
If \(x\in\partial I_\varepsilon^{(n)}\), then by continuity of \(s_n\),
\[
s_n(x)=\varepsilon.
\]
Hence the smallest eigenvalue of
\[
(x\I_n-\A_n(a))^{\mathsf T}(x\I_n-\A_n(a))
\]
equals \(\varepsilon^2\), so
\[
\det\!\bigl((x\I_n-\A_n(a))^{\mathsf T}(x\I_n-\A_n(a))-\varepsilon^2\I_n\bigr)=0.
\]
By Proposition~\ref{prop:factorization},
\[
\det\!\bigl(x\I_n-(\A_n(a)+\varepsilon\PiM_n)\bigr)\,
\det\!\bigl(x\I_n-(\A_n(a)-\varepsilon\PiM_n)\bigr)=0,
\]
which is equivalent to the stated alternative.
\end{proof}

\begin{remark}
Proposition~\ref{prop:factorization} turns the real-axis boundary equation
\[
s_n(x)=\varepsilon
\]
into the vanishing of a product of two ordinary characteristic polynomials.  This
provides an exact algebraic reformulation of the real-slice boundary problem.
\end{remark}

\section{Pentadiagonal determinant formulas for the real-axis singular-value polynomial}
\label{app:pentadiagonal}

Let
\[
\A_n(a)=\tridiag(a,0,1)\in \R^{n\times n},
\qquad 0<a<1,
\]
and define
\[
K_n(x):=(x\I_n-\A_n(a))^{\mathsf T}(x\I_n-\A_n(a)),
\qquad
M_n(x,t):=K_n(x)-t\I_n.
\]
Set
\[
\rho:=x^2+a^2-t,\qquad
p:=x^2+a^2+1-t,\qquad
q:=-(1+a)x,\qquad
s:=a,\qquad
\widetilde\rho:=x^2+1-t.
\]
Then \(M_n(x,t)\) is the symmetric pentadiagonal matrix
\[
M_n(x,t)=
\begin{pmatrix}
\rho & q & s \\
q & p & q & s \\
s & q & p & \ddots \\
& \ddots & \ddots & \ddots & q \\
& & s & q & \widetilde\rho
\end{pmatrix},
\]
with the obvious zeros outside the bandwidth \(2\).

Define the one-corner auxiliary matrix
\[
\mathbf E_n(\rho):=
\begin{pmatrix}
\rho & q & s \\
q & p & q & s \\
s & q & p & \ddots \\
& \ddots & \ddots & \ddots & q \\
& & s & q & p
\end{pmatrix},
\qquad
e_n(\rho):=\det \mathbf E_n(\rho),
\]
and let \(u_j\) denote the \(j\)-th Euclidean basis vector.

\begin{proposition}[Rank-one reduction to a one-corner family]
\label{prop:rank-one-reduction}
For every \(n\ge2\),
\[
M_n(x,t)=\mathbf E_n(\rho)+(\widetilde\rho-p)\,u_nu_n^{\mathsf T}
       =\mathbf E_n(\rho)-a^2u_nu_n^{\mathsf T}.
\]
Consequently,
\[
\det M_n(x,t)=e_n(\rho)-a^2 e_{n-1}(\rho).
\]
\end{proposition}

\begin{proof}
The matrices \(M_n(x,t)\) and \(\mathbf E_n(\rho)\) differ only in the \((n,n)\) entry,
and
\[
\widetilde\rho-p=(x^2+1-t)-(x^2+a^2+1-t)=-a^2.
\]
Hence
\[
M_n(x,t)=\mathbf E_n(\rho)-a^2u_nu_n^{\mathsf T}.
\]
Since only the \((n,n)\) entry is changed, cofactor expansion in that entry gives
\[
\det M_n(x,t)=\det \mathbf E_n(\rho)-a^2\,\operatorname{cof}_{nn}(\mathbf E_n(\rho)).
\]
Removing the last row and last column from \(\mathbf E_n(\rho)\) leaves exactly
\(\mathbf E_{n-1}(\rho)\), so
\[
\operatorname{cof}_{nn}(\mathbf E_n(\rho))=e_{n-1}(\rho),
\]
which proves the claim.
\end{proof}

\begin{proposition}[Five-term recurrence for the one-corner determinant]
\label{prop:e-recurrence}
Assume \(q\neq0\) (equivalently, \(x\neq0\)). Then for \(n\ge7\),
\[
e_n(\rho)
=
(p-s)e_{n-1}(\rho)
+(ps-q^2)\bigl(e_{n-2}(\rho)-s\,e_{n-3}(\rho)\bigr)
+s^3(s-p)e_{n-4}(\rho)
+s^5e_{n-5}(\rho),
\]
with initial values
\begin{align*}
e_1(\rho)&=\rho,\\
e_2(\rho)&=p\rho-q^2,\\
e_3(\rho)&=p^2\rho-q^2(\rho-2s)-p(q^2+s^2),\\
e_4(\rho)&=p^3\rho-p^2(q^2+s^2)-p\!\left(2q^2(\rho-s)+\rho s^2\right)
          +q^4+2q^2s(\rho-s)+s^4,\\
e_5(\rho)&=p^4\rho+q^4(\rho-4s)+\rho s^4+2q^2s^2(-\rho+s)-p^3(q^2+s^2)\\
&\qquad
          +p\!\left(2q^4+4q^2\rho s+s^4\right)
          +p^2\!\left(-2\rho s^2+q^2(-3\rho+2s)\right).
\end{align*}
\end{proposition}

\begin{proof}
A direct Laplace expansion along the last row of \(\mathbf E_n(\rho)\) expresses
\(e_n(\rho)\) in terms of \(e_{n-1}(\rho)\) and two adjacent border minors.  Expanding
those border minors again along their last rows yields a closed finite system of linear
recurrences.  Eliminating the auxiliary minors produces the displayed five-term
recurrence, which is the standard continuant relation for one-corner pentadiagonal
Toeplitz determinants.  The initial values follow from direct determinant evaluation
for \(n=1,2,3,4,5\).
\end{proof}

\begin{corollary}[Closed form in the generic case]\label{cor:generic-closed-form}
Assume
\[
q\neq0,
\qquad
q^2\notin\left\{4s(p-2s),\ \frac{(p+2s)^2}{4}\right\}.
\]
Define
\[
\alpha:=\sqrt{(p+2s)^2-4q^2},
\]
\[
\beta_1:=\sqrt{2(p-2s)(p+2s-\alpha)-4q^2},
\qquad
\beta_2:=\sqrt{2(p-2s)(p+2s+\alpha)-4q^2},
\]
and
\[
\mu_1:=\frac{p-2s-\alpha-\beta_1}{4},\qquad
\mu_2:=\frac{p-2s-\alpha+\beta_1}{4},
\]
\[
\mu_3:=\frac{p-2s+\alpha-\beta_2}{4},\qquad
\mu_4:=\frac{p-2s+\alpha+\beta_2}{4},\qquad
\mu_5:=s.
\]
Then \(\mu_1,\dots,\mu_5\) are the distinct roots of the characteristic polynomial
\[
\chi(z):=
z^5-(p-s)z^4-(ps-q^2)(z^3-sz^2)-s^3(s-p)z-s^5,
\]
and therefore
\[
e_n(\rho)=\sum_{j=1}^5 \kappa_j \mu_j^n,
\qquad n\ge1,
\]
where the coefficients \(\kappa_1,\dots,\kappa_5\) are uniquely determined by
\[
\sum_{j=1}^5 \kappa_j \mu_j^m=e_m(\rho),
\qquad m=1,2,3,4,5.
\]
Consequently,
\[
\det M_n(x,t)
=
\sum_{j=1}^5 \kappa_j(\mu_j-a^2)\mu_j^{\,n-1},
\qquad n\ge1.
\]
\end{corollary}

\begin{proof}
The recurrence in Proposition~\ref{prop:e-recurrence} has characteristic polynomial
\(\chi\).  In the stated generic regime, the roots are distinct, so the standard
characteristic-roots ansatz yields
\[
e_n(\rho)=\sum_{j=1}^5 \kappa_j\mu_j^n.
\]
The coefficients are fixed uniquely by matching the first five values
\(e_1(\rho),\dots,e_5(\rho)\).  The determinant identity then follows from
Proposition~\ref{prop:rank-one-reduction}:
\[
\det M_n(x,t)=e_n(\rho)-a^2e_{n-1}(\rho).
\]
\end{proof}

\begin{proposition}[Degenerate line \(x=0\)]\label{prop:x0-degenerate}
If \(x=0\), then \(q=0\) and the recurrence simplifies to
\[
e_n(\rho)=p\,e_{n-1}(\rho)-p s^2 e_{n-3}(\rho)+s^4 e_{n-4}(\rho),
\qquad n\ge6,
\]
with initial values
\[
e_1(\rho)=\rho,\qquad
e_2(\rho)=p\rho,\qquad
e_3(\rho)=p(p\rho-s^2),\qquad
e_4(\rho)=(p^2-s^2)(p\rho-s^2).
\]
If moreover \(p^2\neq4s^2\), then
\[
\nu_1=-s,\qquad \nu_2=s,\qquad
\nu_3=\frac{p-\sqrt{p^2-4s^2}}{2},\qquad
\nu_4=\frac{p+\sqrt{p^2-4s^2}}{2}
\]
are distinct and
\[
e_n(\rho)=\sum_{j=1}^4 \widetilde\kappa_j \nu_j^n,
\qquad n\ge1,
\]
where the coefficients \(\widetilde\kappa_j\) are determined by
\[
\sum_{j=1}^4 \widetilde\kappa_j \nu_j^m=e_m(\rho),
\qquad m=1,2,3,4.
\]
Hence
\[
\det M_n(0,t)
=
\sum_{j=1}^4 \widetilde\kappa_j(\nu_j-a^2)\nu_j^{\,n-1},
\qquad n\ge1.
\]
\end{proposition}

\begin{proof}
Set \(q=0\) in Proposition~\ref{prop:e-recurrence}.  The recurrence factor reduces,
after cancellation of the redundant factor \((z-s)\), to the displayed fourth-order
relation.  The initial values come from direct determinant evaluation.  The closed
form then follows from the characteristic-roots representation of the reduced
recurrence, and the last identity again follows from
Proposition~\ref{prop:rank-one-reduction}.
\end{proof}

\begin{proposition}[Alternative symmetric-corner reduction]
\label{prop:symmetric-corner}
Define the symmetric-corner auxiliary matrix
\[
\widetilde{\mathbf D}_n(\widetilde\rho):=
\begin{pmatrix}
\widetilde\rho & q & s \\
q & p & q & s \\
s & q & p & \ddots \\
& \ddots & \ddots & \ddots & q \\
& & s & q & \widetilde\rho
\end{pmatrix},
\qquad
\widetilde d_n(x,t):=\det \widetilde{\mathbf D}_n(\widetilde\rho).
\]
Then
\[
M_n(x,t)=\widetilde{\mathbf D}_n(\widetilde\rho)-(1-a^2)u_1u_1^{\mathsf T}.
\]
If \(\widetilde{\mathbf D}_n(\widetilde\rho)\) is invertible, then
\[
\det M_n(x,t)
=
\widetilde d_n(x,t)\Bigl(1-(1-a^2)m_n(x,t)\Bigr),
\]
where
\[
m_n(x,t):=u_1^{\mathsf T}\widetilde{\mathbf D}_n(\widetilde\rho)^{-1}u_1.
\]
\end{proposition}

\begin{proof}
The matrices \(M_n(x,t)\) and \(\widetilde{\mathbf D}_n(\widetilde\rho)\) differ only in the
\((1,1)\) entry, and
\[
\rho-\widetilde\rho=(x^2+a^2-t)-(x^2+1-t)=-(1-a^2).
\]
Hence
\[
M_n(x,t)=\widetilde{\mathbf D}_n(\widetilde\rho)-(1-a^2)u_1u_1^{\mathsf T}.
\]
The matrix determinant lemma gives
\[
\det(\widetilde{\mathbf D}_n+uv^{\mathsf T})
=
\det(\widetilde{\mathbf D}_n)\bigl(1+v^{\mathsf T}\widetilde{\mathbf D}_n^{-1}u\bigr),
\]
and applying it with \(u=-(1-a^2)u_1\) and \(v=u_1\) yields the claim.
\end{proof}

\begin{remark}
In the nongeneric repeated-root cases
\[
q^2\in\left\{4s(p-2s),\ \frac{(p+2s)^2}{4}\right\},
\]
or, in the line \(x=0\), when \(p^2=4s^2\), the same procedure yields the usual
polynomial-times-exponential terms \(n^\ell\mu^n\) associated with repeated roots.
\end{remark}

\section{Even-size central reduction at \texorpdfstring{\(x=0\)}{x=0}}
\label{app:even-central}

\begin{lemma}[Even-size central reduction]\label{lem:even-central}
Let \(n=2m\ge2\), \(0<a<1\), and let
\[
\A_{2m}(a)=\tridiag(a,0,1).
\]
Let \(J_m\in\R^{m\times m}\) be the nilpotent shift
\[
(J_m)_{j,j+1}=1,\qquad j=1,\dots,m-1,
\]
and define
\[
L_m:=a\I_m+J_m,\qquad U_m:=\I_m+aJ_m^{\mathsf T}.
\]
Let \(\Pi_{2m}\in\R^{2m\times2m}\) be the odd-even permutation matrix corresponding to
the reordered basis
\[
e_1,e_3,\dots,e_{2m-1},e_2,e_4,\dots,e_{2m}.
\]
Then
\[
\Pi_{2m}^{\mathsf T}\A_{2m}(a)\Pi_{2m}
=
\begin{bmatrix}
0 & U_m\\
L_m & 0
\end{bmatrix}.
\]
Consequently,
\[
s_{2m}(0)=\smin\!\bigl(\A_{2m}(a)\bigr)=\smin(L_m)=\smin(a\I_m+J_m).
\]
Moreover,
\[
s_{2m+2}(0)\le s_{2m}(0).
\]
\end{lemma}

\begin{proof}
The block form is immediate from the action of \(\A_{2m}(a)\) on odd and even
coordinates: odd indices are mapped only to even indices, and even indices only to
odd indices.  In the odd-even reordered basis, the upper-right block is
\(U_m=\I_m+aJ_m^{\mathsf T}\) and the lower-left block is \(L_m=a\I_m+J_m\).

Therefore
\[
\Pi_{2m}^{\mathsf T}\A_{2m}(a)^{\mathsf T}\A_{2m}(a)\Pi_{2m}
=
\begin{bmatrix}
L_m^{\mathsf T}L_m & 0\\
0 & U_m^{\mathsf T}U_m
\end{bmatrix},
\]
so
\[
\smin\!\bigl(\A_{2m}(a)\bigr)
=
\min\{\smin(L_m),\smin(U_m)\}.
\]

Take any \(x\in\C^m\), and define \(y\in\C^m\) by
\[
y_j:=x_{m+1-j},\qquad j=1,\dots,m.
\]
Then \(\|y\|_2=\|x\|_2\).  With the convention \(y_{m+1}:=0\), we have
\[
\|U_mx\|_2^2=\sum_{j=1}^m |y_j+a y_{j+1}|^2,
\qquad
\|L_my\|_2^2=\sum_{j=1}^m |a y_j+y_{j+1}|^2.
\]
Subtracting,
\begin{align*}
\|U_mx\|_2^2-\|L_my\|_2^2
&=
\sum_{j=1}^m\Bigl(|y_j+a y_{j+1}|^2-|a y_j+y_{j+1}|^2\Bigr) \\
&=
(1-a^2)\sum_{j=1}^m\bigl(|y_j|^2-|y_{j+1}|^2\bigr) \\
&=
(1-a^2)|y_1|^2\ge0.
\end{align*}
Hence \(\|U_mx\|_2\ge \|L_my\|_2\) for every \(x\), and since \(x\mapsto y\) is an
isometry of \(\C^m\), it follows that
\[
\smin(U_m)\ge \smin(L_m).
\]
Thus
\[
s_{2m}(0)=\smin\!\bigl(\A_{2m}(a)\bigr)=\smin(L_m).
\]

For the monotonicity in \(m\), note that \(L_m^{\mathsf T}L_m\) is the leading principal
\(m\times m\) submatrix of \(L_{m+1}^{\mathsf T}L_{m+1}\).  By Cauchy interlacing,
\[
\lambda_{\min}\!\bigl(L_{m+1}^{\mathsf T}L_{m+1}\bigr)
\le
\lambda_{\min}\!\bigl(L_m^{\mathsf T}L_m\bigr).
\]
Taking square roots yields
\[
\smin(L_{m+1})\le \smin(L_m),
\]
and therefore
\[
s_{2m+2}(0)\le s_{2m}(0).
\]
\end{proof}

\end{document}
